%!TEX root = ../template.tex
%%%%%%%%%%%%%%%%%%%%%%%%%%%%%%%%%%%%%%%%%%%%%%%%%%%%%%%%%%%%%%%%%%%%
%% chapter3.tex
%% NOVA thesis document file
%%
%% Chapter with Work Proposal
%%%%%%%%%%%%%%%%%%%%%%%%%%%%%%%%%%%%%%%%%%%%%%%%%%%%%%%%%%%%%%%%%%%%

\typeout{NT FILE chapter3.tex}

%\makeatletter
%\newcommand{\ntifpkgloaded}{%
%  \@ifpackageloaded%
%}
%\makeatother

%Proposta de trabalho

\chapter{Work Proposal}
\label{cha:workProposal}

This project begins with the prototype design and the development of the arduino's code and of an interface app. Afterwards, data from different movements is collected and the \gls{ML} 
models are developed. The last goal is to implement a \gls{ML} model into a server connected via WiFi and a Tiny \gls{ML} model into the app giving real-time feedback and analysis of their movements and fatigue 
monitoring.

Basketball is one of the most popular sports worldwide and it is constantly becoming more complex and more competitive and so does the related technology alongside it. Gathering data from the players is crucial 
to improve the efficiency of their practices and their game. As of today, coaches rely in video analysis to help their players with tactical movements and technical skills but this method has its 
limitations such as blocked views or difficult angles for the recording device especially in a five vs five scenario. This is a strong suit of wearable sensors since they can capture data
continuously and with the right design, be unnoticeable for the user and other players.

For basketball in particular, the movements consist in a lot of short sprints, quick changes of pace and direction, lateral movement and high jumps (vertical and horizontal) and its high tense moments are more common 
at the end of games (fourth quarter), specially at the last minutes.

\newpage

\section{Project Description}
\label{sec:project_description}

In this project, for each foot, an Arduino and an \gls{EMG} sensor will be used. The Arduinos have an integrated \gls{IMU} each that capture acceleration and orientation and they can communicate through \gls{BLE}.
Besides the \gls{IMU}, the \gls{EMG} sensors will collect data from the muscles to provide further insight for the movement recognition and injury report.

\begin{figure}[h]
  \centering
  \includegraphics[width=0.9\linewidth]{System_Overview.jpg}
  \caption{System Overview}
  \label{fig:systemOverview}
\end{figure}

\subsection{Prototype}
\label{sec:prototype}

The Arduino chosen was the Nano 33 \gls{BLE} Rev 2 for its size, 9 axis \gls{IMU}, \gls{BLE} connection and versatility. The choice for the \gls{EMG} sensor was the BITalino Electromyography designed for surface \gls{EMG}. This 
sensors were soldered to the Arduinos for power and communication.

\begin{figure}[ht]
    \centering
    \begin{subfigure}[b]{0.45\textwidth}
        \centering
        \includegraphics[width=\textwidth]{deviceTop.jpg}
        \caption{Top side - Arduino}
        \label{fig:portoTop}
    \end{subfigure}
    \hfill 
    \begin{subfigure}[b]{0.45\textwidth}
        \centering
        \includegraphics[width=\textwidth]{deviceBottom.jpg}
        \caption{Bottom side - \gls{EMG} sensor}
        \label{fig:portoBot}
    \end{subfigure}
  
    \caption{Prototype Photos}
    \label{fig:prototypePic}
\end{figure}

\newpage

\subsubsection{Arduino Specifications}
\label{ssec:arduinoSpecs}

This Arduino runs at 3.3V and the \gls{IMU} sensors are the BMI270 (3-axis accelerometer + 3-axis gyroscope) and the BMM150 (3-axis Magnetometer) sensors. It has a 
width of 18mm and length of 45mm weighing only 5g~\cite{arduinoDoc}. The Arduino IDE will be used for the Arduinos' programming.


\begin{figure}[h]
  \centering
  \includegraphics[width=0.8\linewidth]{arduino-pinout.jpg}
  \caption{Arduino Nano \gls{BLE} Rev 2 - Pinout~\cite{arduinoDoc}}
  \label{fig:arduinoPinout}
\end{figure}

\subsubsection{BITalino Electromyography}
\label{ssec:bitalinoEMG}

The sensor has a gain of 1009 and an input of 2.0-3.5V making the output of 3.3V of the Arduino ideal for it. As seen in ~\autoref{fig:emg}, the emg has an input channel for the electrode cables on the left side 
and 4 connection points on the right side:

\begin{itemize}
  \item \gls{VOUT} - Connected through the white wire to the respective input pin of the Arduino
  \item \gls{VCC} - The red wire is connecting this point to the 3.3V pin.
  \item \gls{VCC}/2 - From the same pin that the \gls{VCC} point is connected to, we have a small voltage divider with an Ampop for stabilization to halve the tension (blue wire).
  \item \gls{GND} - The black wire connects the Arduino's ground to this point.
\end{itemize}

\begin{figure}[h]
  \centering
  \includegraphics[width=0.4\linewidth]{emg.jpg}
  \caption{BITalino \gls{EMG} sensor~\cite{bitalino}}
  \label{fig:emg}
\end{figure}

 A circuit was designed to obtain a stable reference for the \gls{EMG} sensor, implemented using a resistive divider followed by an operational amplifier. As shown in figure~\autoref{fig:divisorTensao}, two identical resistors are 
 connected between the 3.3V supply from the arduino and ground resulting on the node in the middle settling at half the voltage. The op-amp's purpose is to stabilize the circuit by lowering the impedance at the output.

\begin{figure}[h]
  \centering
  \includegraphics[width=0.8\linewidth]{divisorTensao.jpg}
  \caption{Halving the \gls{VCC}}
  \label{fig:divisorTensao}
\end{figure}

The analog sensor data are converted into digital values ranged between 0 and $2^10 - 1$ and for practical reasons we need to convert it back to voltage as show in the following expression:

\begin{equation}
  \gls{EMG}(V) = \frac{(\frac{ADC}{2^10}-\frac{1}{2}*VCC)}{G}
  \label{eq:EMGconvertionEQ}
\end{equation}

\begin{equation}
  \gls{EMG}(mV) = \gls{EMG}(V) * 1000
  \label{eq:EMGconvertionEQmV}
\end{equation}

Where ADC is the digital value from the sensor, VCC is the 3.3V provided by the arduino, G is the sensor gain which is 1009 and EMG(V) and EMG(mV) is the signal in Volt and millivolt respectively.

\subsection{Data Acquisition}
\label{sec:data_acquisition}

A fundamental step in this project is the acquisition of reliable and representative data, as the performance of subsequent analysis methods directly depend on the data quality. The \gls{EMG} sensors will be attached to the 
Gastrocnemius Lateralis (see ~\ref{fig:emgPosition}) to capture the muscle activation patterns in basketball movements.

This process will be carried out under a controlled environment ensuring that the movements being made are correctly categorized. This labeling of activities is essential for the machine learning algorithm.
For a richer dataset the players will be asked to do the movements with different directions, intensities, speed and other variations that can be applied depending on the motion.

Special attention will be given to minimizing signal noise and motion artifacts during acquisition, using appropriate electrode placement and filtering techniques. This ensures that the raw signals preserve as much physiological 
information as possible, thereby improving the robustness of later processing and analysis stages.

%Falta imagem do dispositivo colocado


\subsection{Data Processing}
\label{sec:data_processing}

First the arduino converts the digital value of the \gls{EMG} to mV and afterwards it calculates six different features before sending them to the app: 

\begin{itemize}
  \item Root Mean Square - \gls{RMS}
  \item Mean Absolute Value - \gls{MAV}
  \item Waveform Length - \gls{WL}
  \item Zero Crossing - \gls{ZC}
  \item Slope Sign Changes - \gls{SSC}
  \item Variance - \gls{VAR}
\end{itemize}


Following the feature calculations, the arduinos send the data to the android device that forwards it to the database where it gets stored and to the server that predicts which movement 
the user is doing. With the stored data, two \gls{ML} models are trained and then used for basketball movements recognition.

